\documentclass[12pt]{article}
\usepackage[utf8]{inputenc}
\usepackage[T1]{fontenc}
\usepackage[margin=1in]{geometry}
\usepackage{setspace}
\usepackage{xcolor}
\usepackage[colorlinks=true, urlcolor=blue!60!black, linkcolor=blue!60!black]{hyperref}
\setlength{\parindent}{0em}
\setlength{\parskip}{0.45em}
\renewcommand{\baselinestretch}{1.0}
\begin{document}
\begin{center}
{\Large \textit{United States v. Cruikshank}, 92 U.S. 542 (1876)\par}
\vspace{0.3em}
{\normalsize Case Brief --- One-Page Edition\par}
\end{center}

\textbf{Memory Jogger.} The First and Second Amendments bind only Congress, and the Fourteenth reaches only state action; the Colfax Massacre convictions fell (\href{https://supreme.justia.com/cases/federal/us/92/542/}{Justia}).

\textbf{Facts.} In 1873 a white mob massacred Black freedmen guarding the Colfax, Louisiana courthouse. Federal prosecutors charged the attackers under the Enforcement Act of 1870 with conspiring against victims' federal rights.

\textbf{Procedural History.} Three men were convicted; the circuit judges divided on arresting judgment and certified the question.

\textbf{Votes.} Arrest of judgment affirmed, 9--0 (Waite for eight; Clifford concurring on pleading grounds).

\textbf{Holding.} No count stated a federal offense: assembly and arms rights restrict only Congress, the Fourteenth Amendment bars only state action, the voting counts omitted racial motive, and catch-all counts were fatally vague.

\textbf{Reasoning.} Under dual citizenship and \textit{Barron}, pre-existing rights look to state protection; the Fourteenth Amendment guards only against States. The Court ``may suspect that race was the cause of the hostility,'' but it was not averred, and indictments must ``descend to particulars.'' \emph{Separately:} Clifford rested solely on the indictment's lack of particularity.

\textbf{Examples.} \emph{Same-Side:} a purely private assault with no federal nexus. \emph{Opposite-Side:} intimidating Black voters in a federal election. \emph{Fence-Sitter:} a State's systematic refusal to prosecute racial terror.

\textbf{Critique.} Condemned for betraying Reconstruction and immunizing white supremacist terror.

\textbf{Key Quotation.} ``This is not a right granted by the Constitution. Neither is it in any manner dependent upon that instrument for its existence. The second amendment declares that it shall not be infringed, but this, as has been seen, means no more than that it shall not be infringed by Congress.'' --- Chief Justice Waite, Opinion of the Court.

\end{document}
